%!TEX root=../../main.tex

The antecedents of convex neural networks --- both finite and
infinite-dimensional --- arose from Bayesian analyses.
This line of work begins with \citet{neal1996priors}, who showed that
infinite-width, two-layer neural networks converge to
Gaussian process models given appropriate priors on the network weights.
\citet{williams1996computing} then gave an exact form for the covariance
function of these processes.
Recently, many works have extended the Gaussian process characterization to
large-width limits of deep neural networks
\citep{lee2019deep, matthews2018deep}, convolutional networks
\citep{novak2019convolutional, alonso2019convolutional}, and recurrent
architectures \citep{yang2019recurrent}.
However, these models rely on a probabilistic approach that views the network
weights as random variables and does not account for training dynamics.
Convexity of the final prediction model is a by-product of the analysis, rather
than the goal.

Another convexification of infinite-width neural networks is given by
the neural tangent kernel (NTK) \citep{jacot2018ntk}.
The NTK innovates on Bayesian analyses by extending the Gaussian process
characterization to include the behavior of the network during training;
this is achieved by showing that the effects of optimization vanish in the
infinite-width limit, causing the network to converge to a Gaussian process
with a kernel specified by the gradients at initialization.
NTK dynamics have been the target of intense study by the machine learning
community and extensions have been developed for convolutions
\citep{arora2019ntk}, graph neural networks \citep{du2019graph} and general
architectures \citep{yang2020tensor}, among other models.

A major criticism of the NTK approach is that it does not capture feature
learning.
Since the effects of training vanish as the network becomes arbitrarily wide,
NTK models remain arbitrarily close to their initializations in the so-called
``lazy training regime'' \citep{chizat2019lazy}.
The mean-field, or ``rich training regime'', is an alternative large-width
limit for two-layer neural networks where the model converges to a signed
measure over an uncountably infinite number of basis functions
\citep{bach2017breaking, mei2018mean}.
In this formulation, the prediction function is an affine map with respect to
the basis expansion, meaning training mean-field networks is a fully convex,
although infinite-dimensional, problem.
The mean-field approach originates with \citet{bengio2006convex} and
\citet{rosset2007regularization}, who use convex boosting sub-problems to find
optimal networks where the number of units is chosen automatically by a
sparsity inducing norm.
More recent work has focused on the convergence of gradient flows for the
measure optimization problem \citep{chizat2018convergence} and extensions
of the mean-field regime to deeper networks
\citep{yang2020tensor, sirignano2022mean, chizat2024speed}.

Although the mean-field regime has been key to demonstrating provable
feature-learning in neural networks, infinite-width problems are impractical to
solve and do not yield direct insights into finite training problems.
Convex reformulations solve both of these problems by providing finite-size,
convex versions of neural networks.
Key steps towards convex reformulations were achieved by
\citet{bartan2019relaxations}, who gave a convex relaxation of convolutional
networks, and \citet{ergen2020convex}, who derived a convex, semi-infinite dual
problem for two-layer ReLU models.
\citet{pilanci2020convexnn} gave the first proper convex reformulation by
computing the associated bi-dual and simplifying it using the piece-wise
linear structure of ReLU activations.
The resulting problem is equivalent to a finite discretization of the
mean-field limit for two-layer ReLU networks.

The space of convex reformulations has steadily expanded since
\citet{pilanci2020convexnn}.
In addition to the works mentioned at the start of this chapter, convex
reformulations have been developed for deep parallel networks
\citep{ergen2021beyond, wang2023parallel}, threshold activations
\citep{ergen2023threshold}, and quantized networks \citep{bartan2021quantized}.
Other papers have explored the connections between convex reformulations and
the NTK \citep{dwaraknath2023ntk} and established approximation bounds
for lower-dimensional relaxations \citep{kim2024convex}.
\citet{lacotte2020hidden} use convex reformulations to show that spurious local
minima vanish if the network is sufficiently wide, while \citet{wang2022hidden}
show that local minima correspond to sub-sampled convex reformulations.

Finally, we briefly mention input convex networks
\citep{amos2017input, sivaprasad2021curious}.
Input convex networks are neural networks for which the predictions are a
convex function of the input features, rather than the model parameters.
This approach is completely distinct from convex reformulations and does not
address the challenge of the non-convex training problem.
